\section{Irreducibility of the ordinary point on every open subgroup}
\label{sec:spectrum-open-irreducibility}

The rank-three automorphy endpoint in
Section~\ref{sec:spectrum-deformation} needs irreducibility on every
open subgroup contained in $G_K$. We prove this for the ordinary
integral type-I points in that section. The essential restriction on
the residual character is visible in the proof.

Write $G=G_{\mathbb Q}$, $H=\ker\eta=G_K$, and $I_\ell$ for inertia at
the chosen valuation above $\ell$. Let $\mathcal O\subset A\subset
\overline{\mathbb Q}_\ell$ be the finite DVR of the ordinary point. Use
$k_A$ for its residue field and write
$\Psi=\chi_\ell^e\psi_{\mathrm{Gal}}$ in the coefficient units, with
$\psi_{\mathrm{Gal}}$ of Dirichlet level $N$ and $e\ge1$. On $H$ the
matrix part of the point is an ordinary representation, denoted by
$R$. In the marked integral basis its reduction has diagonal entries
$\overline\Psi,1,\overline\Psi^{-1}$ on $H$ (equivalently one may
use just this residual trace identity below). At the decomposition
group at $\ell$ intersected with $H$, the ordinary flag gives a
basis in which the matrices are upper triangular. On inertia the
diagonal characters in that basis are $\Psi,1,\Psi^{-1}$.
The arithmetic input \texttt{dat\_facts} gives
\begin{equation}
 I_\ell\subset H,\qquad
 \exists x\in I_\ell:\quad \Psi(x)^2-1\in\mathcal O^\times .
 \label{eq:open-residue-separation}
\end{equation}
In particular this difference also has nonzero image in $k_A$.
The hypothesis that $R$ acts irreducibly on $H$ is precisely type I.

\subsection{A three-dimensional permutation argument}

Let $H$ act irreducibly on $V=L^3$ over a field $L$, let $U$ be a
normal subgroup of $H$, and suppose some $u\in U$ has three distinct
eigenvalues in $L$. Suppose $R(u)$ and $R(x)$ are upper triangular
in a common basis for some $x\in H$. If $U$ acts reducibly then there
is a $g\in H$ for which
\begin{equation}
               \operatorname{tr} R(gx^j)=0\qquad(j\ge0).
 \label{eq:open-vanishing-coset}
\end{equation}
To see this, diagonalize $R(u)$ by an upper triangular change of
basis. For each diagonal position its eigenvector is obtained by
choosing entry one at that position, zero below it, and solving
successively in the rows above it. The differences of diagonal
entries that must be divided by are all nonzero. Thus $R(x)$
remains upper triangular. Denote the three resulting eigenlines
by $L_0,L_1,L_2$.

Any $U$-invariant subspace is a sum of some of these lines: the
projection to the $i$th line is the polynomial in $R(u)$ given by
$\prod_{a\ne i}(R(u)-\lambda_a)/(\lambda_i-\lambda_a)$.
Take a nonzero proper $U$-invariant subspace $W$. If it is a line,
then a $U$-invariant eigenline has been found. If it is a plane,
irreducibility on $H$ gives $h\in H$ with $hW\ne W$.
Normality makes $hW$ another $U$-invariant plane. Their intersection
is a $U$-invariant line. Every translate by $H$ of a $U$-invariant
line is again $U$-invariant, hence one of the three eigenspaces.
The span of its orbit is nonzero and $H$-invariant. Irreducibility
implies that the orbit consists of all three lines. Thus the action
on the three lines gives a transitive subgroup of $S_3$.
Its order is divisible by three, so it contains a three-cycle
by Cauchy's theorem. Choose $g$ realizing this cycle.

Every $R(h)$ is now a monomial matrix. The permutation of an upper
triangular invertible monomial matrix fixes each index: if it sends
column $i$ to row $p(i)$, then $p(i)\le i$ for every $i$, and
$\sum_i p(i)=\sum_i i$. In particular $R(x)$ is diagonal.
$R(gx^j)$ therefore has the same permutation as $R(g)$ and all its
diagonal entries vanish. This proves \eqref{eq:open-vanishing-coset}.

\subsection{Inertia supplies both matrices}

Let $V_0$ be an arbitrary open subgroup of $G$ contained in $H$,
and set $U=\bigcap_{g\in G}gV_0g^{-1}$. Openness in the profinite
group gives finite index, so the normal core $U$ also has finite
index. It suffices to prove irreducibility on $U$, as every
$V_0$-invariant subspace is $U$-invariant. Fix the element $x$ in
\eqref{eq:open-residue-separation}. The cyclotomic inertia theorem
over $\mathbb Q_\ell$ gives $\sigma\in I_\ell$ with
$\chi_\ell(\sigma)=1+\ell$. Put
\[
    m=[G:U]\, |(\mathbb Z/N\mathbb Z)^\times|,\qquad u=\sigma^m .
\]
Here $m>0$, $u\in U\cap I_\ell$, and $\psi_{\mathrm{Gal}}(u)=1$,
since its order divides the order of the Dirichlet unit group.
It follows that $\Psi(u)^2=(1+\ell)^{2em}\ne1$ in characteristic
zero. In the ordinary basis the three eigenvalues of $R(u)$ are
therefore distinct. Both $u$ and $x$ belong to inertia, so the
ordinary flag triangularizes both matrices.

If $U$ acted reducibly, \eqref{eq:open-vanishing-coset} would produce
$g\in H$. Compute its vanishing traces in the integral basis, and
reduce modulo the maximal ideal of $A$. Put $a=\overline\Psi(g)$,
$b=\overline\Psi(x)$ in $k_A^\times$. The residual trace identity
gives
\[
                 ab^j+1+(ab^j)^{-1}=0\quad(j\ge0).
\]
Then each of $a,ab,ab^2$ is a root of $X^2+X+1$.
Since $b^2\ne1$ by \eqref{eq:open-residue-separation}, $b\ne1$
and the three units are pairwise distinct. A quadratic polynomial
over a field has at most two distinct roots. This contradiction proves
irreducibility on $U$ and hence on $V_0$. The proof works for each
open $V_0\subseteq H$ and supplies the required quantifier.

\paragraph{Source and inputs.}
This is \texttt{matIrr\_open\_gd}. The argument above derives the open
irreducibility conclusion from the residual trace, the ordinary flag,
\eqref{eq:open-residue-separation}, and cyclotomic surjectivity on
inertia. For the given good deformation data the residual and ordinary
conditions belong to \texttt{OrdPt}; \texttt{dat\_facts} supplies the
displayed separation by the character and unramifiedness conditions
of the datum. The characteristic-zero and residual eigenvalue
conditions have different uses. The power $u$ distinguishes the
three eigenlines in characteristic zero, and $x$ distinguishes three
residual values after applying the Frobenius-independent trace test.
\footnote{\href{https://github.com/Deicyde/riemann-hypothesis/blob/40e56bab026552e089d7afcfeff08d2defb9c005/src/Development/AutomorphicArithmetic/AutomorphicVarieties/CurveDVRSpecialization/MatIrrOpenGd.lean}{The invariant subspaces, trace obstruction and arithmetic specialization},
\href{https://github.com/Deicyde/riemann-hypothesis/blob/40e56bab026552e089d7afcfeff08d2defb9c005/src/Development/AutomorphicArithmetic/AutomorphicVarieties/UnitaryThreeGroups/DatFacts.lean}{inertia and residual separation}.}
